% ============================================================
%  chapters/chap5_contribution_technique.tex
%  Contribution technique : la plateforme FasoTALN
% ============================================================

% ----- Bibliothèque TikZ requise par le diagramme de cas d'utilisation -----
% (à déplacer dans config/preamble.tex si d'autres chapitres en ont besoin)
\usetikzlibrary{shapes.geometric}

% ----- Styles TikZ locaux (valables pour tout ce chapitre) -----
\tikzset{
	Bn/.style={rectangle, rounded corners=2pt, draw=black!55, fill=gray!7,
		minimum height=1.0cm, text centered, font=\small},
	Bh/.style={rectangle, rounded corners=2pt, draw=black!70, fill=gray!20,
		minimum height=1.0cm, text centered, font=\small},
	Bf/.style={rectangle, rounded corners=2pt, draw=black!70, fill=gray!14,
		minimum height=1.1cm, text centered, font=\small},
	Bl/.style={rectangle, rounded corners=3pt, draw=black!45, fill=gray!3,
		inner sep=8pt},
	fl/.style={-{Stealth[length=5pt, width=4pt]}, thick, draw=black!60},
	fld/.style={-{Stealth[length=5pt, width=4pt]}, thick, draw=black!60, dashed},
	flb/.style={{Stealth[length=5pt, width=4pt]}-{Stealth[length=5pt, width=4pt]},
		thick, draw=black!60},
}

% ----- Emplacement réservé pour une capture d'écran non encore produite -----
% Permet au document de compiler et aux références croisées de fonctionner
% avant que la capture ne soit insérée. Remplacer par \figurefull{...} ensuite.
\newcommand{\captureareserver}[3]{%
	\begin{figure}[H]
		\centering
		\fbox{\parbox[c][3.2cm][c]{0.72\textwidth}{\centering\footnotesize
				\textit{Capture d'écran à insérer~:} \texttt{#1}}}
		\caption{#2}
		\label{fig:#3}
	\end{figure}
}

\chapter{Contribution technique~: la plateforme FasoTALN}
\label{chap:contribution-technique}

\section{Introduction}

Les chapitres~\ref{chap:conception} et~\ref{chap:realisation} ont établi une
contribution scientifique~: un module G2P pour le \langueafr{Mooré}, le
\langueafr{Dioula} et le \langueafr{Bambara}, et la mesure de l'apport d'une
représentation phonémique dans un scénario de transfert cross-lingual zero-shot.
Cette contribution existe sous la forme de deux modèles publiés sur HuggingFace et
de notebooks d'expérimentation. Sous cette forme, elle reste inaccessible à qui ne
sait pas exécuter un notebook, et invisible pour la communauté qu'elle concerne au
premier chef~: les étudiants, enseignants et praticiens burkinabè du TALN.

Ce chapitre présente la réponse apportée à ce constat~: \textbf{FasoTALN}, un portail
scientifique, pédagogique et technologique consacré au traitement automatique des
langues nationales du Burkina Faso, qui remplit simultanément deux fonctions. Il
\textbf{industrialise} le pipeline de recherche~--- les deux modèles deviennent un
service d'inférence interrogeable en ligne, sans installation~--- et il
\textbf{contextualise} ce pipeline en le replaçant dans un parcours d'apprentissage
progressif qui va des enjeux linguistiques des langues nationales jusqu'aux
perspectives de recherche ouvertes.

La frontière avec les chapitres précédents est nette et volontaire~: le
chapitre~\ref{chap:realisation} répond à la question \og{}le signal phonémique
améliore-t-il la classification cross-lingue~?\fg{}, ce chapitre-ci répond à la
question \og{}comment ce résultat devient-il un système utilisable et un objet
d'enseignement~?\fg{}. Aucune métrique expérimentale n'est donc reprise ici~; en
revanche, les contraintes d'implémentation qui conditionnent la fidélité du service
au modèle entraîné y sont détaillées, car elles constituent le c\oe{}ur du travail
d'ingénierie.

Le chapitre suit la démarche classique de conduite de projet logiciel~: analyse et
spécification des besoins (section~\ref{sec:ct-besoins}), conception globale et
détaillée (section~\ref{sec:ct-conception}), puis réalisation
(section~\ref{sec:ct-realisation}).

\section{Analyse et spécification des besoins}
\label{sec:ct-besoins}

\subsection{Analyse des besoins}

\subsubsection{Identification des acteurs}

La plateforme est un portail public de consultation~: elle ne comporte ni
authentification, ni compte utilisateur, ni espace privé. Deux acteurs humains et
trois acteurs secondaires (systèmes externes) interagissent avec elle~:

\begin{itemize}
	\item \textbf{Le visiteur-apprenant}~: acteur principal. Il s'agit d'un étudiant,
	d'un enseignant, d'un chercheur ou d'un curieux qui découvre le TALN appliqué aux
	langues nationales. Il consulte le contenu pédagogique, exécute les démonstrations
	interactives et interroge l'assistant conversationnel. Il n'a aucun droit
	d'écriture sur le contenu.

	\item \textbf{Le mainteneur scientifique}~: l'auteur du projet ou un contributeur
	ultérieur. Il enrichit la base de connaissance éditoriale (fichiers JSON
	versionnés), intègre de nouveaux résultats expérimentaux et peut déclencher
	manuellement la mise à jour de la veille d'actualité. Il agit par le dépôt Git et
	par un point d'entrée d'administration technique, non par une interface graphique
	d'administration.

	\item \textbf{L'ordonnanceur}~: acteur secondaire non humain. Il déclenche
	périodiquement (toutes les 24~heures) la mise à jour de la veille d'actualité,
	sans intervention humaine.

	\item \textbf{Les services externes}~: l'API Mistral (assistant conversationnel et
	agent de veille) et l'API CITADEL (traduction français~$\rightarrow$~mooré). Ce
	sont des acteurs secondaires sollicités par le système.
\end{itemize}

\subsubsection{Analyse des besoins fonctionnels}

La plateforme doit répondre aux besoins suivants~:

\begin{itemize}
	\item proposer un \textbf{parcours d'apprentissage ordonné} qui conduit le
	visiteur, étape par étape, des langues nationales du Burkina Faso jusqu'aux
	perspectives de recherche, en passant par les défis du domaine, les approches
	existantes et les ressources disponibles~;
	\item permettre la \textbf{transcription phonémique interactive} d'un texte saisi
	par l'utilisateur en \langueafr{Mooré}, \langueafr{Dioula} ou \langueafr{Bambara}~;
	\item permettre l'\textbf{exécution complète du pipeline de classification}
	(texte~$\rightarrow$~IPA~$\rightarrow$~séquence hybride~$\rightarrow$~classe
	thématique) en exposant chaque état intermédiaire, et non seulement le résultat
	final~;
	\item \textbf{générer une phrase de test en mooré} à partir d'une phrase française,
	afin qu'un visiteur ne lisant aucune des langues cibles puisse néanmoins
	expérimenter le pipeline~;
	\item \textbf{servir une base de connaissance éditoriale sourcée} couvrant les
	langues, le glossaire, les défis, les approches, les ressources, l'écosystème
	institutionnel et les perspectives~;
	\item offrir une \textbf{recherche transversale} sur l'ensemble du contenu~;
	\item mettre à disposition un \textbf{assistant conversationnel} capable de
	répondre aux questions du visiteur en s'appuyant sur cette base de connaissance~;
	\item entretenir une \textbf{veille d'actualité} sur le TALN africain, mise à jour
	automatiquement et systématiquement sourcée.
\end{itemize}

\subsubsection{Analyse des besoins non fonctionnels}

Six exigences non fonctionnelles structurent la solution. Elles ne sont pas
génériques~: chacune découle d'une contrainte propre au contexte du projet.

\begin{itemize}
	\item \textbf{Fiabilité scientifique du contenu.} Le portail s'adresse à un public
	qui découvre le domaine et n'a pas les moyens de vérifier ce qu'il lit. Toute
	affirmation factuelle externe doit donc être accompagnée d'une source vérifiable,
	et aucune valeur numérique ne doit être approximée. Cette exigence s'applique
	aussi bien au contenu rédigé qu'au contenu produit automatiquement, ce qui
	conditionne toute l'architecture de l'assistant conversationnel et de la veille
	d'actualité (section~\ref{sec:ct-realisation}).

	\item \textbf{Dégradation gracieuse.} La plateforme dépend de services externes
	(Mistral, CITADEL) dont la disponibilité n'est pas garantie et dont les clés
	peuvent être absentes d'un déploiement donné. L'indisponibilité d'un service ne
	doit jamais empêcher l'accès au reste du portail~: elle doit se traduire par un
	message ciblé sur la seule fonctionnalité concernée.

	\item \textbf{Évolutivité éditoriale.} Le projet est conçu pour vivre plusieurs
	années et accueillir de nouvelles langues, de nouvelles tâches et de nouveaux
	résultats. L'ajout d'une page de contenu ne doit pas exiger la modification de la
	barre de navigation, du plan du site, du fil d'Ariane, de la navigation
	séquentielle ni du moteur de recherche.

	\item \textbf{Maîtrise des coûts d'exploitation.} L'assistant conversationnel
	repose sur une API facturée à l'usage, sur un portail public sans
	authentification. Une limitation de débit est donc nécessaire pour éviter qu'un
	usage abusif ne rende le service insoutenable.

	\item \textbf{Sobriété d'exécution.} La cible de déploiement est un hébergement
	gratuit, sans accélérateur matériel. Les deux modèles (environ 1,5~Go cumulés)
	doivent donc s'exécuter sur processeur seul, dans une empreinte mémoire
	compatible avec cette contrainte.

	\item \textbf{Accessibilité et adaptabilité.} Une part importante du public cible
	consulte le web depuis un téléphone mobile et une connexion contrainte. L'interface
	est donc conçue selon une approche \textit{mobile-first}, et chaque image
	illustrative porte une description textuelle alternative.
\end{itemize}

\subsection{Spécification des besoins}

\subsubsection{Diagramme de cas d'utilisation}

La figure~\ref{fig:ct-usecase} présente les cas d'utilisation de la plateforme. Le
visiteur-apprenant dispose de huit cas d'utilisation, tous accessibles sans
authentification. Deux d'entre eux~--- l'exécution du pipeline de classification et
la génération d'une phrase de test~--- sont liés par une relation
\code{<<extend>>}~: la traduction française~$\rightarrow$~mooré est une commodité
optionnelle offerte au sein de la page du pipeline, et non une étape obligatoire.
La transcription phonémique, en revanche, est incluse (\code{<<include>>}) dans
l'exécution du pipeline, puisque la séquence hybride ne peut être construite sans
la sortie du module G2P.

Le mainteneur scientifique n'accède à aucune interface d'administration~: il met à
jour le contenu par le dépôt versionné et peut déclencher la veille. L'ordonnanceur
partage ce dernier cas d'utilisation, ce qui traduit un choix de conception
important~: le rafraîchissement automatique et le rafraîchissement manuel empruntent
exactement le même chemin d'exécution, et donc les mêmes garde-fous.

% -------------------------------------------------------------------
% FIGURE — Diagramme de cas d'utilisation
% -------------------------------------------------------------------
\begin{figure}[H]
	\centering
	\begin{tikzpicture}[x=1cm, y=1cm]

		\tikzset{
			uc/.style={ellipse, draw=black!60, fill=gray!5, align=center,
				font=\scriptsize, text width=2.9cm, inner sep=2.5pt, minimum height=0.95cm},
			lien/.style={draw=black!60, thick},
			rel/.style={-{Stealth[length=4pt, width=3pt]}, thick, draw=black!60, dashed},
		}

		% ---- Cas d'utilisation du visiteur (colonne 1) ----
		\node[uc] (u1) at (0,   0.00) {Suivre le parcours d'apprentissage};
		\node[uc] (u2) at (0,  -1.45) {Consulter la base de connaissance};
		\node[uc] (u3) at (0,  -2.90) {Rechercher dans le portail};
		\node[uc] (u4) at (0,  -4.35) {Consulter les résultats expérimentaux};
		\node[uc] (u5) at (0,  -5.80) {Transcrire un texte en IPA (G2P)};
		\node[uc] (u6) at (0,  -7.25) {Exécuter le pipeline de classification};
		\node[uc] (u7) at (0,  -8.70) {Interroger l'assistant};
		\node[uc] (u8) at (0, -10.15) {Consulter la veille d'actualité};

		% ---- Cas d'utilisation liés / administration (colonne 2) ----
		\node[uc] (u9)  at (5.2,  -7.25) {Générer une phrase de test en mooré};
		\node[uc] (u10) at (5.2,  -9.60) {Rafraîchir la veille d'actualité};
		\node[uc] (u11) at (5.2, -11.00) {Mettre à jour la base éditoriale};

		% ---- Frontière du système ----
		\begin{scope}[on background layer]
			\node[draw=black!55, rounded corners=2pt, fill=gray!2, inner sep=12pt,
			fit=(u1)(u8)(u9)(u11)] (sys) {};
		\end{scope}
		\node[font=\footnotesize\bfseries, anchor=north west] at ($(sys.north west)+(0.25,-0.12)$)
		{Plateforme FasoTALN};

		% ---- Acteur : visiteur-apprenant ----
		\begin{scope}[shift={(-4.2,-5.10)}]
			\draw[lien] (0,0.42) circle (0.19);
			\draw[lien] (0,0.23) -- (0,-0.32);
			\draw[lien] (-0.30,0.02) -- (0.30,0.02);
			\draw[lien] (0,-0.32) -- (-0.25,-0.78);
			\draw[lien] (0,-0.32) -- (0.25,-0.78);
			\node[font=\scriptsize, align=center, anchor=north] at (0,-0.90) {Visiteur-\\apprenant};
		\end{scope}
		\coordinate (V) at (-3.9,-5.10);

		% ---- Acteur : ordonnanceur ----
		\begin{scope}[shift={(9.6,-8.40)}]
			\draw[lien] (0,0.42) circle (0.19);
			\draw[lien] (0,0.23) -- (0,-0.32);
			\draw[lien] (-0.30,0.02) -- (0.30,0.02);
			\draw[lien] (0,-0.32) -- (-0.25,-0.78);
			\draw[lien] (0,-0.32) -- (0.25,-0.78);
			\node[font=\scriptsize, align=center, anchor=north] at (0,-0.90) {Ordonnanceur};
		\end{scope}
		\coordinate (O) at (9.3,-8.40);

		% ---- Acteur : mainteneur scientifique ----
		\begin{scope}[shift={(9.6,-10.60)}]
			\draw[lien] (0,0.42) circle (0.19);
			\draw[lien] (0,0.23) -- (0,-0.32);
			\draw[lien] (-0.30,0.02) -- (0.30,0.02);
			\draw[lien] (0,-0.32) -- (-0.25,-0.78);
			\draw[lien] (0,-0.32) -- (0.25,-0.78);
			\node[font=\scriptsize, align=center, anchor=north] at (0,-0.90) {Mainteneur\\scientifique};
		\end{scope}
		\coordinate (M) at (9.3,-10.60);

		% ---- Associations du visiteur ----
		\foreach \n in {u1,u2,u3,u4,u5,u6,u7,u8}{
			\draw[lien] (V) -- (\n.west);
		}

		% ---- Associations des acteurs secondaires ----
		\draw[lien] (O) -- (u10.east);
		\draw[lien] (M) -- (u10.east);
		\draw[lien] (M) -- (u11.east);

		% ---- Relations entre cas d'utilisation ----
		\draw[rel] (u1.east) to[bend left=55] node[right, font=\scriptsize, xshift=1pt]
		{$\ll$include$\gg$} (u2.east);
		\draw[rel] (u6.north east) to[bend left=55] node[right, font=\scriptsize, xshift=1pt]
		{$\ll$include$\gg$} (u5.south east);
		\draw[rel] (u9.west) -- node[above, font=\scriptsize] {$\ll$extend$\gg$} (u6.east);

	\end{tikzpicture}
	\caption{Diagramme de cas d'utilisation de la plateforme FasoTALN.}
	\label{fig:ct-usecase}
\end{figure}

\subsubsection{Diagrammes de séquence}

Trois scénarios méritent d'être détaillés, parce qu'ils portent chacun une décision
de conception non triviale.

\paragraph{Scénario 1~— Exécution du pipeline de classification.}
La figure~\ref{fig:ct-seq-pipeline} décrit le scénario nominal. Le point notable est
la boucle de génération~: le module G2P ne transcrit pas la phrase en une passe, mais
mot par mot, chaque mot étant préfixé par le code de langue avant d'être soumis au
modèle. Ce découpage n'est pas un détail d'optimisation, c'est une contrainte de
fidélité au format vu par le modèle pendant son entraînement, détaillée en
section~\ref{sec:ct-realisation}. Le service renvoie ensuite au client non seulement la classe
prédite, mais aussi la transcription IPA intermédiaire et la séquence hybride
textuelle complète, afin que l'interface puisse rendre visible chaque étape du
raisonnement.

\figurefull{figures/uml_sequence_pipeline.png}{0.86}
{Diagramme de séquence~: exécution du pipeline de classification.}{ct-seq-pipeline}

\paragraph{Scénario 2~— Interrogation de l'assistant conversationnel.}
La figure~\ref{fig:ct-seq-chat} présente le scénario alternatif de l'assistant. Trois
issues sont possibles, matérialisées par un fragment \code{alt}~: la réponse est
contenue dans la base de connaissance injectée au démarrage~; la question porte sur un
fait récent absent de cette base, et l'agent recourt alors à une recherche web
restreinte à une liste de domaines de confiance, en citant sa source~; ou bien rien
de fiable n'est trouvé, et l'agent le déclare explicitement plutôt que de produire
une réponse plausible mais invérifiable. Le contrôle de débit intervient en amont de
tout appel externe, de sorte qu'une requête rejetée ne consomme aucun crédit d'API.

\figurefull{figures/uml_sequence_chat.png}{0.97}
{Diagramme de séquence~: assistant conversationnel, ancrage documentaire et repli
	de recherche web.}{ct-seq-chat}

\paragraph{Scénario 3~— Mise à jour de la veille d'actualité.}
La figure~\ref{fig:ct-seq-news} décrit le scénario le plus sensible de la plateforme,
puisque du contenu y est publié sans relecture humaine. Deux mécanismes de rejet y
apparaissent~: si la réponse du modèle n'est pas un document JSON exploitable, la
mise à jour est abandonnée et le contenu existant est conservé intact~; et chaque
élément est ensuite filtré individuellement selon le domaine de sa source. Le
principe directeur est qu'aucune défaillance ne doit pouvoir dégrader l'état
existant~: dans le pire des cas, la page n'est pas rafraîchie.

\figurefull{figures/uml_sequence_news.png}{0.82}
{Diagramme de séquence~: mise à jour de la veille d'actualité et filtrage des
	sources.}{ct-seq-news}

\section{Conception}
\label{sec:ct-conception}

\subsection{Conception globale}

\subsubsection{Architecture physique}

L'architecture physique retenue, présentée en figure~\ref{fig:ct-archi-physique},
s'écarte volontairement du modèle trois-tiers classique. Elle ne comporte que
\textbf{deux tiers}~: un client léger (le navigateur du visiteur, exécutant une
application monopage) et un serveur d'application unique qui héberge à la fois
l'interface compilée, l'API et les deux modèles chargés en mémoire.

L'absence du troisième tiers~--- le serveur de données~--- est un choix assumé et non
une omission. La plateforme ne collecte aucune donnée utilisateur~: ni compte, ni
historique de conversation persistant, ni enregistrement audio, ni journal des textes
soumis aux démonstrations. Il n'y a donc rien à stocker durablement, et par
conséquent aucune base de données à administrer, à sauvegarder ni à sécuriser. Ce
choix a trois conséquences favorables~: le service est \textit{sans état} et peut être
redémarré ou redéployé sans procédure de migration~; la surface d'exposition en
matière de données personnelles est nulle~; et l'ensemble tient dans un conteneur
unique, ce qui rend le déploiement sur un hébergement gratuit réaliste.

Le contenu éditorial, qui pourrait justifier une base de données dans une application
comparable, est traité comme du \textbf{code}~: il vit dans des fichiers versionnés
avec le reste du projet, ce qui lui confère un historique de modification, une
possibilité de relecture avant intégration et une reproductibilité que ne donnerait
pas une base administrée par une interface.

% -------------------------------------------------------------------
% FIGURE — Architecture physique
% -------------------------------------------------------------------
\begin{figure}[H]
	\centering
	\begin{tikzpicture}[x=1cm, y=1cm]

		% ---- Client ----
		\node[Bl, minimum width=4.0cm, minimum height=3.6cm] (cli) at (-5.4, 0) {};
		\node[font=\small\bfseries, anchor=north] at (-5.4, 2.05) {Client léger};
		\node[Bn, text width=3.2cm] at (-5.4, 0.45)
		{Navigateur\\[2pt]\footnotesize Application monopage};
		\node[font=\footnotesize, text width=3.4cm, align=center] at (-5.4, -1.15)
		{Rendu, navigation\\et état local};

		% ---- Serveur d'application ----
		\node[Bl, minimum width=5.8cm, minimum height=5.4cm] (srv) at (0.6, -0.5) {};
		\node[font=\small\bfseries, anchor=north, text width=6.4cm, align=center]
		at (0.6, 3.45) {Serveur d'application\\(conteneur unique)};
		\node[Bn, text width=4.6cm] at (0.6, 1.4) {Interface compilée\\[2pt]\footnotesize Fichiers statiques};
		\node[Bn, text width=4.6cm] at (0.6, -0.1) {API REST\\[2pt]\footnotesize FastAPI};
		\node[Bh, text width=4.6cm] at (0.6, -1.7) {Modèles en mémoire\\[2pt]\footnotesize ByT5 G2P~+~AfroXLM-R};

		% ---- Services externes ----
		\node[Bl, minimum width=4.0cm, minimum height=3.6cm] (ext) at (6.6, 0) {};
		\node[font=\small\bfseries, anchor=north] at (6.6, 2.05) {Services externes};
		\node[Bn, text width=3.2cm] at (6.6, 0.55) {\footnotesize API Mistral};
		\node[Bn, text width=3.2cm] at (6.6, -0.75) {\footnotesize API CITADEL};

		% ---- Absence de serveur de données ----
		\node[rectangle, rounded corners=3pt, draw=black!35, dashed, fill=none,
		text width=6.4cm, align=center, font=\footnotesize, inner sep=6pt]
		(nodb) at (0.6, -4.3)
		{\textit{Aucun serveur de données~:}\\aucune donnée utilisateur persistée};

		% ---- Flèches ----
		\draw[flb] (-3.4, 0) -- (-2.3, 0);
		\node[font=\footnotesize, anchor=south] at (-2.85, 0.12) {HTTP};
		\draw[flb] (3.5, 0) -- (4.6, 0);
		\node[font=\footnotesize, anchor=south] at (4.05, 0.12) {HTTPS};
		\draw[black!35, dashed] (0.6, -3.2) -- (0.6, -3.85);

	\end{tikzpicture}
	\caption{Architecture physique~: deux tiers, sans serveur de données.}
	\label{fig:ct-archi-physique}
\end{figure}

\subsubsection{Architecture logique}

Les composants sont organisés en trois couches (figure~\ref{fig:ct-archi-logique}),
chacune n'adressant que la couche immédiatement inférieure.

La \textbf{couche présentation} regroupe les pages de contenu, les composants
d'interface réutilisables et les gabarits de mise en page. Elle est pilotée par un
\textbf{descripteur de navigation unique}~: un fichier de configuration qui déclare
les groupes thématiques et l'ordre curriculaire des pages. La barre de navigation, le
plan latéral, le fil d'Ariane, la navigation séquentielle \og{}précédent/suivant\fg{},
le décompte des étapes affiché sur la page d'accueil et l'index du moteur de recherche
en sont tous dérivés. C'est la réponse concrète à l'exigence d'évolutivité éditoriale~:
déclarer une page dans ce descripteur suffit à l'insérer dans l'ensemble des
mécanismes de navigation du site.

La \textbf{couche applicative} expose les points d'entrée de l'API et héberge six
composants métier indépendants~: les deux moteurs d'inférence, le moteur
conversationnel, le moteur de traduction, l'agent de veille et le chargeur de contenu.
Chacun encapsule intégralement sa dépendance externe et signale son indisponibilité
par une erreur typée, ce qui permet à la couche API de la traduire en réponse
\code{503} ciblée~--- mécanisme par lequel l'exigence de dégradation gracieuse est
satisfaite.

La \textbf{couche contenu} est constituée des fichiers JSON éditoriaux. Elle est en
lecture seule pour l'ensemble du système, à une exception près~: le fichier des
actualités, seul artefact écrit à l'exécution, et uniquement par l'agent de veille.

% -------------------------------------------------------------------
% FIGURE — Architecture logique
% -------------------------------------------------------------------
\begin{figure}[H]
	\centering
	\begin{tikzpicture}[x=1cm, y=1cm]

		% ---- Couche présentation ----
		\node[Bl, minimum width=14.2cm, minimum height=2.9cm] at (0, 3.4) {};
		\node[font=\small\bfseries, anchor=west] at (-6.9, 4.45) {Couche présentation};
		\node[Bn, text width=2.8cm] at (-4.95, 2.9) {\scriptsize Pages de contenu};
		\node[Bn, text width=2.8cm] at (-1.65, 2.9) {\scriptsize Gabarit\\documentaire};
		\node[Bh, text width=2.8cm] at (1.65, 2.9) {\scriptsize Descripteur\\de navigation};
		\node[Bn, text width=2.8cm] at (4.95, 2.9) {\scriptsize Assistant (widget)};

		% ---- Couche applicative ----
		\node[Bl, minimum width=14.2cm, minimum height=4.1cm] at (0, -0.8) {};
		\node[font=\small\bfseries, anchor=west] at (-6.9, 0.95) {Couche applicative};
		\node[Bn, text width=12.6cm] at (0, 0.0) {\scriptsize Points d'entrée de l'API (FastAPI)};
		\node[Bh, text width=2.8cm] at (-4.95, -1.7) {\scriptsize Moteur G2P};
		\node[Bh, text width=2.8cm] at (-1.65, -1.7) {\scriptsize Classifieur hybride};
		\node[Bn, text width=2.8cm] at (1.65, -1.7) {\scriptsize Moteur\\conversationnel};
		\node[Bn, text width=2.8cm] at (4.95, -1.7) {\scriptsize Agent de veille};

		% ---- Couche contenu ----
		\node[Bl, minimum width=14.2cm, minimum height=2.0cm] at (0, -4.35) {};
		\node[font=\small\bfseries, anchor=west] at (-6.9, -3.65) {Couche contenu};
		\node[Bf, text width=12.6cm] at (0, -4.6)
		{\scriptsize Base de connaissance éditoriale~--- fichiers JSON versionnés (lecture seule)};

		% ---- Flèches inter-couches ----
		\draw[fl] (0, 1.95) -- (0, 1.25);
		\node[font=\footnotesize, anchor=west] at (0.15, 1.6) {HTTP / JSON};
		\draw[fl] (0, -2.85) -- (0, -3.35);

	\end{tikzpicture}
	\caption{Architecture logique en trois couches.}
	\label{fig:ct-archi-logique}
\end{figure}

\subsection{Conception détaillée}

\subsubsection{Aspect statique}

\paragraph{Structure des composants métier.}
Les composants de la couche applicative se répartissent en deux familles, présentées
séparément pour la lisibilité~: les composants d'\textbf{inférence}
(figure~\ref{fig:ct-classes-inf}), qui exécutent les modèles localement, et les
composants de \textbf{service et de contenu} (figure~\ref{fig:ct-classes-srv}), qui
encapsulent les dépendances externes et l'accès à la base éditoriale.

\figurefull{figures/uml_classes_inference.png}{0.62}
{Diagramme de classes~: composants d'inférence.}{ct-classes-inf}

Deux principes structurent la première famille. D'une part, les deux moteurs
d'inférence chargent leur modèle \textbf{une seule fois}, à l'instanciation~: le coût
de chargement (téléchargement initial, désérialisation des poids) est ainsi payé au
démarrage du service et non à chaque requête. D'autre part, les deux moteurs
s'ignorent mutuellement~: c'est la couche API qui compose le pipeline en enchaînant le
moteur G2P puis le classifieur. Cette absence de couplage permet d'exposer la
transcription phonémique comme fonctionnalité autonome, et de faire évoluer chaque
modèle indépendamment. Les méthodes privées \code{word\_to\_ipa} et
\code{build\_hybrid\_inputs} concentrent à elles seules les contraintes de fidélité au
format d'entraînement, détaillées en section~\ref{sec:ct-realisation}.

\figurefull{figures/uml_classes_services.png}{0.74}
{Diagramme de classes~: composants de service et de contenu.}{ct-classes-srv}

Dans la seconde famille, deux dépendances méritent d'être relevées. Le moteur
conversationnel dépend du chargeur de contenu, dont il sérialise la totalité au
démarrage pour constituer son ancrage documentaire~; et il importe de l'agent de
veille la liste des domaines de confiance, de sorte que les deux composants ne
puissent pas diverger sur ce qui constitue une source acceptable.

\paragraph{Modèle de contenu éditorial.}
En l'absence de base de données, le modèle de données du système est constitué de
dix documents JSON, décrits dans le tableau~\ref{tab:ct-contenu}. Chaque document
correspond à une catégorie servie par un même point d'entrée paramétré. Deux formes
coexistent~: la plupart des catégories sont des \textit{listes} d'entités homogènes,
tandis que deux d'entre elles~--- le panorama linguistique et l'écosystème~--- sont
des \textit{objets} structurés en sections, leur contenu étant de nature narrative
plutôt qu'énumérative.

\begin{table}[H]
	\centering
	\caption{Base de connaissance éditoriale~: catégories et rôles}
	\label{tab:ct-contenu}
	\begin{tabularx}{\textwidth}{@{} L{3.4cm} L{1.9cm} X @{}}
		\toprule
		\textbf{Catégorie} & \textbf{Forme} & \textbf{Contenu} \\
		\midrule
		\code{languages}          & liste  & Fiches des langues nationales et du bambara \\
		\code{languages\_context} & objet  & Panorama linguistique, familles de langues, cadre institutionnel \\
		\code{glossaire}          & liste  & Vocabulaire du TALN par catégorie \\
		\code{challenges}         & liste  & Défis scientifiques du domaine, avec sources \\
		\code{approaches}         & liste  & Approches méthodologiques actuelles, avec sources \\
		\code{resources}          & liste  & Corpus, jeux de données, modèles et outils recensés \\
		\code{perspectives}       & liste  & Perspectives de recherche ouvertes par la contribution \\
		\code{ecosysteme}         & objet  & Institutions, réseaux panafricains, parcours de spécialisation \\
		\code{results}            & objet  & Résultats expérimentaux définitifs de la contribution \\
		\code{news}               & liste  & Veille d'actualité (seul fichier écrit à l'exécution) \\
		\bottomrule
	\end{tabularx}
\end{table}

Trois de ces catégories~--- les défis, les approches et l'écosystème~--- portent un
champ \code{sources} normalisé (titre, auteurs, année, lieu de publication, lien),
rendu à l'écran par un composant dédié. Ce champ est le support technique de
l'exigence de fiabilité scientifique~: il rend structurellement visible toute
affirmation non sourcée, et interdit qu'une donnée externe soit ajoutée sans
référence vérifiable.

\paragraph{Points d'entrée de l'API.}
Le tableau~\ref{tab:ct-api} récapitule l'interface de programmation exposée. Elle est
volontairement réduite~: neuf points d'entrée, sans authentification, sans état de
session.

\begin{table}[H]
	\centering
	\caption{Points d'entrée de l'API}
	\label{tab:ct-api}
	\begin{tabularx}{\textwidth}{@{} L{1.3cm} L{4.5cm} X @{}}
		\toprule
		\textbf{Verbe} & \textbf{Route} & \textbf{Rôle} \\
		\midrule
		GET  & \code{/api/health}                & État du service et modèles chargés \\
		POST & \code{/api/g2p}                   & Transcription phonémique d'un texte \\
		POST & \code{/api/pipeline/classify}     & Pipeline complet~: IPA, séquence hybride, classe et probabilités \\
		POST & \code{/api/translate}             & Traduction français~$\rightarrow$~mooré \\
		POST & \code{/api/chat}                  & Assistant conversationnel (débit limité) \\
		GET  & \code{/api/knowledge/\{cat\}}     & Contenu éditorial d'une catégorie \\
		GET  & \code{/api/results}               & Résultats expérimentaux \\
		GET  & \code{/api/news}                  & Veille d'actualité \\
		POST & \code{/api/news/refresh}          & Déclenchement d'une mise à jour de la veille \\
		\bottomrule
	\end{tabularx}
\end{table}

Les deux points d'entrée d'inférence n'acceptent que les trois codes de langue pour
lesquels un modèle existe (\code{mos}, \code{dyu}, \code{bam}). Le fulfuldé et le
gourmantché, bien que documentés dans le contenu éditorial, sont explicitement
signalés comme non couverts par les démonstrations~--- une manière d'exposer l'état
réel de la couverture linguistique plutôt que de laisser croire à une généralité que
les modèles n'ont pas.

\subsubsection{Aspect dynamique}

\paragraph{Exécution du pipeline.}
La figure~\ref{fig:ct-act-pipeline} détaille le diagramme d'activités du pipeline de
classification. Il fait apparaître la décomposition en mots de la phrase d'entrée, la
boucle de transcription, la reconstruction manuelle de la séquence hybride puis la
classification. Le plafonnement du nombre de mots traités y figure explicitement~: la
génération étant séquentielle et exécutée sur processeur, le temps de réponse croît
linéairement avec la longueur du texte, ce qui impose une borne pour préserver la
réactivité du service.

\figurefull{figures/uml_activite_pipeline.png}{0.42}
{Diagramme d'activités~: exécution du pipeline de classification.}{ct-act-pipeline}

\paragraph{Mise à jour de la veille.}
La figure~\ref{fig:ct-act-news} présente le diagramme d'activités de l'agent de
veille. Les deux points de décision correspondent aux deux garde-fous décrits plus
loin (section~\ref{sec:ct-realisation})~: validité syntaxique de la réponse, puis
appartenance de chaque source à la liste des domaines de confiance. La fusion avec
l'historique existant est réalisée par déduplication sur le lien de la source, ce qui
garantit qu'un même article annoncé plusieurs fois n'apparaît qu'une fois dans la
liste publiée.

\figurefull{figures/uml_activite_news.png}{0.78}
{Diagramme d'activités~: mise à jour de la veille d'actualité.}{ct-act-news}

\section{Réalisation}
\label{sec:ct-realisation}

\subsection{Environnement de travail}

\subsubsection{Environnement matériel}

Le développement de la plateforme et l'exécution locale des modèles ont été réalisés
sur un poste de travail personnel. L'entraînement des modèles, décrit au
chapitre~\ref{chap:conception}, a quant à lui été mené sur des accélérateurs
graphiques infonuagiques et n'est pas repris ici. La cible de déploiement diffère
sensiblement du poste de développement~: il s'agit d'un conteneur hébergé sans
accélérateur graphique, ce qui impose que l'inférence soit viable sur processeur seul.

\begin{table}[H]
	\centering
	\caption{Environnements matériels de développement et de déploiement}
	\label{tab:ct-materiel}
	\begin{tabularx}{\textwidth}{@{} L{4.2cm} X X @{}}
		\toprule
		 & \textbf{Développement} & \textbf{Déploiement} \\
		\midrule
		Type          & Poste de travail personnel & Conteneur hébergé (HuggingFace Spaces) \\
		Calcul        & Processeur (inférence), GPU infonuagique pour l'entraînement & Processeur uniquement \\
		Modèles chargés & ByT5 G2P + AfroXLM-R ($\approx$~1,5~Go) & Identiques \\
		Persistance   & Aucune & Aucune (système de fichiers éphémère) \\
		\bottomrule
	\end{tabularx}
\end{table}

\subsubsection{Environnement logiciel}

Les choix technologiques (tableau~\ref{tab:ct-logiciel}) répondent à trois critères~:
continuité avec l'environnement d'expérimentation, aptitude à servir une application
web publique, et compatibilité avec un hébergement gratuit.

\begin{table}[H]
	\centering
	\caption{Environnement logiciel de la plateforme}
	\label{tab:ct-logiciel}
	\begin{tabularx}{\textwidth}{@{} L{3.1cm} L{3.4cm} X @{}}
		\toprule
		\textbf{Catégorie} & \textbf{Outil} & \textbf{Rôle dans le projet} \\
		\midrule
		Serveur applicatif & FastAPI (Python)      & Exposition de l'API, validation des requêtes, cycle de vie des modèles \\
		Inférence          & PyTorch, Transformers & Chargement et exécution des deux modèles, sur processeur \\
		Interface          & React, Vite           & Application monopage, rendu du contenu et des démonstrations \\
		Mise en forme      & Tailwind CSS          & Système de styles fondé sur des variables de conception \\
		Animation          & Framer Motion         & Transitions et animations des démonstrations \\
		Iconographie       & Lucide React          & Jeu d'icônes vectorielles (aucun émoji dans l'interface) \\
		Services externes  & API Mistral, API CITADEL & Assistant conversationnel, veille, traduction \\
		Conteneurisation   & Docker                & Image unique interface~+~API \\
		Hébergement        & HuggingFace Spaces    & Déploiement public du portail \\
		Versionnement      & Git / GitHub          & Historique du code et du contenu éditorial \\
		\bottomrule
	\end{tabularx}
\end{table}

Le choix de FastAPI plutôt qu'un serveur applicatif généraliste tient à un point
précis~: son mécanisme de cycle de vie permet d'initialiser les modèles au démarrage
du processus et de les conserver en mémoire pour toute la durée de vie du service, ce
qui est la condition d'un temps de réponse acceptable pour une inférence sur
processeur.

\subsection{Travail réalisé}

\subsubsection{Industrialisation du module G2P en service d'inférence}
\label{ssec:ct-g2p}

Cette première tâche consistait à faire passer le module G2P du notebook
d'expérimentation au service exécuté en continu. L'enjeu n'était pas la performance
mais la \textbf{fidélité}~: obtenir en production exactement les transcriptions
observées à l'évaluation. Trois contraintes se sont révélées déterminantes.

\paragraph{Chargement du modèle avec tête de sortie non liée.}
Le modèle G2P a été affiné avec une tête de sortie \code{lm\_head} indépendante des
embeddings d'entrée. Or l'implémentation de référence de l'architecture T5 lie ces
deux jeux de poids par défaut au chargement~: sans instruction contraire, la tête de
sortie entraînée est silencieusement écrasée par les poids d'embedding. Le modèle se
charge alors sans erreur, mais produit une sortie proche du hasard. Le chargement
force donc explicitement la dissociation des poids.

\paragraph{Tokenizer natif plutôt que tokenizer du dépôt.}
Le service instancie le tokenizer depuis le modèle de base et non depuis le dépôt du
modèle affiné. Cette précaution répond à des divergences de segmentation constatées
lors des mises au point, qui se traduisaient par des transcriptions dégradées sans
qu'aucune erreur ne soit levée.

\paragraph{Génération mot à mot avec préfixe de langue.}
Le modèle n'a jamais vu, pendant son entraînement, ni phrase complète ni entrée non
préfixée. Le service reproduit donc strictement le format d'entraînement~: chaque mot
est mis en minuscules, préfixé du code de langue, soumis individuellement, et les
transcriptions obtenues sont réassemblées avec un séparateur explicite. Les
paramètres de génération reprennent ceux de l'évaluation. Une version antérieure de ce
composant introduisait des pénalités de répétition~: elles ont été retirées après
analyse, car elles ne faisaient que masquer les répétitions provoquées par un format
d'entrée incorrect, au lieu de corriger la cause.

Un post-traitement phonologique (nasalisation, affriquées), transposé depuis le dépôt
du modèle, est appliqué au \langueafr{Dioula} et au \langueafr{Bambara}~; le
\langueafr{Mooré} n'en requiert pas. La page de démonstration correspondante
(figure~\ref{fig:ct-capture-g2p}) permet de saisir un texte, de choisir la langue et
d'obtenir la transcription phonémique.

\captureareserver{figures/captures/demo\_g2p.png}
{Page de démonstration de la transcription phonémique (G2P).}{ct-capture-g2p}

\subsubsection{Intégration du classifieur hybride}
\label{ssec:ct-classifieur}

La seconde tâche portait sur l'intégration du classifieur. Elle a mis au jour la
difficulté d'implémentation la plus coûteuse du projet~: la \textbf{construction de la
séquence d'entrée hybride}.

Le format attendu par le modèle est \code{[CLS] texte [SEP] IPA [SEP]}, avec un
séparateur \textbf{unique} entre les deux segments. L'appel standard du tokenizer pour
une paire de segments produit en revanche le format de paire hérité de RoBERTa, qui
insère \textbf{deux} séparateurs consécutifs. Ce format ne correspond pas à celui vu
par le modèle pendant l'affinage, et dégrade fortement les prédictions~--- là encore
sans qu'aucune erreur ne soit signalée.

La séquence est donc construite manuellement~: chaque segment est encodé séparément,
sans tokens spéciaux et tronqué à la moitié de la longueur maximale, puis les
identifiants sont assemblés dans l'ordre voulu avec les séparateurs explicites. La
troncature symétrique garantit qu'un texte long ne peut évincer la transcription
phonémique de la fenêtre du modèle~--- ce qui reviendrait à annuler silencieusement
l'apport que le pipeline cherche précisément à exploiter.

L'interface correspondante (figure~\ref{fig:ct-capture-pipeline}) rend visible chaque
état intermédiaire~: texte saisi, transcription IPA produite, séquence hybride
assemblée, puis classe prédite accompagnée de la distribution de probabilités sur les
cinq catégories. Ce parti pris pédagogique est la raison pour laquelle le service
renvoie la séquence assemblée en clair~: elle n'a aucune utilité fonctionnelle pour le
client, mais elle donne à voir ce qui est habituellement invisible dans une
démonstration de classification.

\captureareserver{figures/captures/pipeline.png}
{Page du pipeline de classification~: états intermédiaires et distribution de
	probabilités.}{ct-capture-pipeline}

\subsubsection{Construction du portail pédagogique}
\label{ssec:ct-portail}

La troisième tâche a consisté à concevoir le portail lui-même~: douze pages de contenu
organisées en cinq groupes thématiques, formant un parcours ordonné qui va des langues
nationales aux perspectives de recherche.

Le principe de réalisation est la \textbf{source de vérité unique} décrite en
conception~: l'ordre curriculaire est déclaré une fois, et six mécanismes distincts en
sont dérivés~--- menu de navigation, plan latéral, fil d'Ariane, navigation
séquentielle, décompte des étapes de la page d'accueil et index de recherche. Ajouter
une page se réduit ainsi à deux gestes~: créer le composant et déclarer son entrée dans
le descripteur.

Toutes les pages de contenu partagent un gabarit de type documentation, qui fournit un
en-tête, un plan latéral, un sommaire de page synchronisé avec le défilement et une
navigation \og{}précédent/suivant\fg{}. Le sommaire est construit par observation
dynamique du contenu affiché, ce qui lui permet de fonctionner identiquement pour les
pages statiques et pour celles dont le contenu est chargé de façon asynchrone depuis
l'API.

Une révision complète de l'identité visuelle a par ailleurs été menée en cours de
projet. La palette initiale, héritée d'un produit antérieur à vocation
culturelle et touristique, évoquait l'artisanat plus que la science~; elle a été
remplacée par une palette claire de type plateforme d'apprentissage, tout en
conservant les noms des variables de conception afin de ne pas invalider les usages
existants dans les composants. Une signature graphique discrète, inspirée des motifs
géométriques du bogolan, a été conservée comme unique élément d'ancrage culturel.

\captureareserver{figures/captures/accueil.png}
{Page d'accueil~: parcours d'apprentissage numéroté.}{ct-capture-accueil}

\captureareserver{figures/captures/page\_contenu.png}
{Gabarit documentaire d'une page de contenu~: plan latéral, sommaire et navigation
	séquentielle.}{ct-capture-contenu}

\subsubsection{Assistant conversationnel ancré}
\label{ssec:ct-chat}

La quatrième tâche répondait à un besoin observé~: un visiteur qui découvre le domaine
ne sait pas toujours dans quelle page chercher sa réponse. Un assistant
conversationnel a donc été intégré, sous forme de bulle flottante accessible depuis
toutes les pages.

La difficulté n'est pas de faire répondre un modèle de langue, mais de l'empêcher de
répondre \textit{n'importe quoi} sur un sujet où les faits vérifiables sont rares et
où une erreur serait reprise de bonne foi. Trois mécanismes y concourent.

\textbf{Ancrage documentaire.} L'intégralité de la base de connaissance éditoriale,
résultats expérimentaux et veille compris, est sérialisée \textbf{une seule fois au
	démarrage} et injectée dans les instructions de l'agent. Le corpus étant de taille
modérée et entièrement maîtrisé, cette approche évite la complexité d'une recherche
sémantique par vectorisation tout en garantissant que l'agent dispose de la totalité
du contenu du site, et non d'un extrait sélectionné par une heuristique de similarité.

\textbf{Périmètre explicite.} L'agent est instruit de refuser poliment toute question
hors du champ du TALN et des langues africaines, plutôt que d'y répondre avec
l'assurance d'un modèle généraliste.

\textbf{Repli de recherche sourcé.} Pour une question portant sur un fait récent absent
de la base, l'agent peut effectuer une recherche web, restreinte aux mêmes domaines de
confiance que l'agent de veille~--- la liste est importée depuis ce dernier, de sorte
que les deux composants ne peuvent pas diverger. La source est alors systématiquement
ajoutée à la réponse. Si rien de fiable n'est trouvé, l'agent déclare son ignorance.

Deux protections complètent le dispositif~: une limitation de débit par adresse IP,
évaluée avant tout appel externe, et un plafonnement de l'historique de conversation
transmis, qui borne la taille des requêtes. Enfin, en l'absence de clé d'API,
l'assistant signale son indisponibilité sans affecter le reste du portail.

\captureareserver{figures/captures/assistant.png}
{Assistant conversationnel~: réponse ancrée sur la base de connaissance.}{ct-capture-chat}

\subsubsection{Veille d'actualité automatisée}
\label{ssec:ct-news}

La cinquième tâche visait à maintenir le portail vivant entre deux contributions
éditoriales, en recensant l'actualité du TALN africain. La contrainte est frontale~:
il s'agit de \textbf{publier du contenu sans relecture humaine} sur un portail dont la
crédibilité scientifique est l'exigence première.

La solution retenue est un \textbf{agent unique}~--- et non un système multi-agents~---
qui effectue en une passe la recherche et la synthèse, au moyen d'un connecteur de
recherche web. Deux garde-fous indépendants encadrent sa production.

Le premier est \textbf{déclaratif}~: des instructions strictes interdisent l'usage de
connaissances internes non vérifiées, imposent que chaque lien soit une URL réellement
retournée par la recherche, exigent une valeur nulle plutôt qu'une date inventée, et
énumèrent une liste de domaines de confiance (archives de prépublications, actes de
conférences, réseaux de recherche africains, institutions de financement, ministères
burkinabè).

Le second est \textbf{impératif}~: le serveur vérifie lui-même, avant toute écriture,
que le domaine de chaque source appartient à cette liste. Tout élément non conforme
est écarté. Ce doublon est délibéré~: un prompt est une consigne, pas une garantie~;
la conformité ne pouvait pas reposer sur la seule obéissance du modèle.

Le comportement en cas d'échec obéit au même principe de prudence~: si la réponse
n'est pas exploitable ou si tous les éléments sont rejetés, le contenu existant est
conservé intact. La liste publiée est fusionnée avec l'historique par déduplication
sur le lien de la source, triée par date décroissante et plafonnée.

Le déclenchement est double~: une boucle interne rafraîchit la veille toutes les
24~heures, et un point d'entrée dédié permet un déclenchement externe. Cette
redondance répond à une limite concrète de l'hébergement gratuit~: un conteneur mis en
veille faute de trafic n'exécute plus sa boucle interne, et un déclencheur externe
périodique devient alors le mécanisme fiable.

\captureareserver{figures/captures/nouvelles.png}
{Page \og{}Nouvelles du jour\fg{}~: actualités avec catégorie, date et lien
	source.}{ct-capture-news}

\subsubsection{Réintroduction maîtrisée de la traduction}

La sixième tâche répond à un obstacle d'usage~: la démonstration du pipeline suppose de
saisir une phrase en mooré, en dioula ou en bambara, ce dont la majorité des visiteurs
est incapable. Une traduction français~$\rightarrow$~mooré a donc été intégrée à la
page du pipeline, afin de produire une phrase de test.

Ce composant est le seul élément réemployé du produit antérieur, et son périmètre a
été délibérément restreint~: uniquement du texte, uniquement vers le mooré, sans
synthèse vocale ni sens inverse, et sans identifiants inscrits dans le code~--- ceux-ci
proviennent exclusivement de l'environnement d'exécution. Les limites du service
externe sous-jacent (absence de prise en charge du dioula et du bambara) sont
affichées explicitement dans l'interface plutôt que masquées.

\subsubsection{Conteneurisation et déploiement}

La dernière tâche a porté sur la mise en ligne. Le déploiement s'appuie sur une image
Docker construite en deux étapes~: la première compile l'interface, la seconde
installe l'environnement Python et n'y recopie que le résultat de la compilation. Ce
découpage évite d'embarquer la chaîne d'outils front-end dans l'image finale.

Deux ajustements ont été nécessaires pour la cible d'hébergement. D'une part, la
variante processeur seul de la bibliothèque de calcul est installée explicitement, ce
qui évite d'embarquer plusieurs gigaoctets de dépendances graphiques inutilisables sur
la cible. D'autre part, le répertoire de cache des modèles est redirigé vers un
emplacement inscriptible, l'hébergeur exécutant le conteneur sous un utilisateur non
privilégié.

En production, l'API sert également l'interface compilée~: les deux composants
partagent la même origine, ce qui supprime la nécessité d'une configuration de partage
de ressources entre origines. Un mécanisme de repli renvoie le document d'entrée de
l'application pour toute route inconnue, de manière à ce que le routage côté client
fonctionne lors d'un accès direct à une adresse profonde.

\subsubsection{Difficultés rencontrées et enseignements}

Trois difficultés ont marqué cette phase de réalisation. Elles partagent une
caractéristique qui mérite d'être soulignée~: \textbf{aucune ne provoquait d'erreur}.
Le service démarrait, répondait, et produisait des résultats plausibles mais faux.

La première est le déliement de la tête de sortie du modèle G2P~: un chargement
apparemment nominal produisait des transcriptions aberrantes. La deuxième est le
double séparateur inséré par l'appel standard du tokenizer pour une paire de
segments~: les prédictions restaient syntaxiquement valides, mais nettement moins
bonnes qu'à l'évaluation. La troisième, plus insidieuse, est l'ajout de pénalités de
répétition à la génération~: introduites pour corriger un symptôme, elles masquaient
la cause réelle, à savoir un format d'entrée non conforme.

L'enseignement méthodologique est le suivant~: pour un composant issu d'un travail
expérimental, \textbf{le notebook d'entraînement et d'évaluation est la spécification},
et la seule validation recevable d'une modification est la comparaison des sorties du
service avec celles observées dans ce notebook. Une modification jugée équivalente par
raisonnement ne l'est pas nécessairement pour le modèle. Cette règle a été inscrite
dans la documentation du projet, afin qu'elle survive à son auteur.

\section{Conclusion}

Ce chapitre a présenté la contribution technique du travail~: la plateforme FasoTALN.
L'analyse des besoins a identifié ses acteurs et ses exigences, au premier rang
desquelles la fiabilité scientifique du contenu et la dégradation gracieuse face à
l'indisponibilité des services externes. La conception a établi une architecture à
deux tiers, sans base de données, organisée en trois couches et pilotée par une source
de vérité unique pour la navigation. La réalisation a détaillé sept tâches, de
l'industrialisation des modèles de recherche au déploiement conteneurisé.

La contribution propre de ce volet est double. D'une part, elle rend accessible sans
installation un pipeline de recherche jusque-là confiné à des notebooks, en exposant
non seulement son résultat mais chacun de ses états intermédiaires, ce qui en fait un
support d'enseignement autant qu'une démonstration. D'autre part, elle documente les
conditions de fidélité d'un service d'inférence à son modèle d'origine~--- un ensemble
de contraintes qui, si elles ne sont pas respectées, dégradent silencieusement les
performances mesurées en laboratoire.
